\section{第~\ref{sec:dofaalt}~章：\nt{DOPA}的其他理论}

每一种扩展图景都有自己对多巴胺的直接测量作为依据（翻转点的离散、对厌恶和新异事件的响应、缓慢上升、对味道替换的响应），但解释它们的公式是理论，其中两种理论之间的实验结果目前相互矛盾。

{\footnotesize
\setlength{\tabcolsep}{3pt}\renewcommand{\arraystretch}{1.15}
\begin{longtable}{>{\raggedright}p{0.5cm}>{\raggedright}p{4.0cm}>{\raggedright}p{5.6cm}>{\raggedright}p{2.3cm}>{\raggedright\arraybackslash}p{1.7cm}}
\toprule
序 & 命题 & 实验与测量内容 & 来源 & 状态 \\
\midrule\endfirsthead
\toprule
序 & 命题 & 实验与测量内容 & 来源 & 状态 \\
\midrule\endhead
1 & 不对称规则~\eqref{eq:asymmetric}收敛到点~\eqref{eq:expectile}；$\kappa=1/2$时为平均值，对概率为$1/2$的一滴奖励有$V=\kappa$（图~\ref{fig:expectiles}） & 点~\eqref{eq:expectile}是期望分位数（expectile），即正负偏差权重不同的最小二乘问题的解；本章例子的推导见本章。 & \cite{NeweyPowell1987}；本章推导 & 理论 \\
2 & 不同\nt{DOPA}神经元的翻转点不同，并按不对称性排序（命题~\ref{prop:expectile}） & 小鼠，腹侧被盖区经光遗传标记的\nt{DOPA}神经元，不同大小的奖励：爆发变为低谷的点因神经元而异；响应不对称性越大的神经元，翻转点越高；由一组神经元的响应可重建奖励分布的形状。这是离散的主要功能，则是假说。 & \cite{Dabney2020} & 已测量 \\
3 & 部分\nt{DOPA}神经元对厌恶事件以爆发响应 & 猴子：有的\nt{DOPA}神经元被奖励预示信号兴奋、被眼部吹气的预示信号抑制，另一些则对二者都兴奋；这两组位于中脑的不同部分。 & \cite{Matsumoto2009, BrombergMartin2010} & 已测量 \\
4 & 误差的模或新异性设定学习速率 & 所引文献中没有\nt{DOPA}神经元的重要性信号改变学习速率的直接实验；综述提出了“注意，重要”信号的作用。 & \cite{BrombergMartin2010} & 假说 \\
5 & 危险维度有单独的导线 & 小鼠：投射到纹状体尾部的\nt{DOPA}神经元对新异的强刺激响应，而不对奖励响应；毁损它们会减弱对新物体的回避。 & \cite{Menegas2018} & 已测量 \\
6 & 最优动作时间$\tau_a^*=\sqrt{C/\bar R}$~\eqref{eq:vigor} & 和式$C/\tau_a+\bar R\tau_a$的最小值；本章推导。在原始模型中，动作速度由等于平均奖励的时间代价决定。 & \cite{Niv2007}；本章推导 & 理论 \\
7 & 多巴胺慢电平携带$\bar R$并设定动作速度（命题~\ref{prop:multiplex}） & 大鼠，伏隔核中的微透析和伏安法：多巴胺水平逐分钟跟随奖励频率和动作速度。在人类中，L-DOPA增强反应时间对平均奖励频率的依赖。慢电平恰好等于$\bar R$，尚未证明。 & \cite{Hamid2016, Beierholm2013vigor} & 假说 \\
8 & 慢电平由两个来源组成：输出端的释放可在脉冲频率不变时改变，但频率本身也有缓慢上升 & 大鼠，同一任务：伏隔核中的释放跟随变化的奖励期望，而腹侧被盖区已鉴定\nt{DOPA}神经元的脉冲频率则不然。虚拟走廊中的小鼠（第~10~行）：已鉴定\nt{DOPA}神经元的脉冲中也有接近奖励时的平缓上升。 & \cite{Mohebi2019, Kim2020} & 已测量 \\
9 & 动物接近远处奖励时多巴胺平缓上升 & 迷宫中的大鼠，纹状体伏安法：随着接近目标，多巴胺浓度在几秒内上升，斜率取决于距离和奖励大小。 & \cite{Howe2013ramp, Hamid2016} & 已测量 \\
10 & 上升是估计修正时的误差$\delta$~\eqref{eq:ramp}；走廊中的跳变产生爆发（图~\ref{fig:ramp}） & 公式和每秒$\delta=0.75\,V$这个数是本章推导。实验：虚拟走廊中的小鼠，瞬间移到离目标更近处；胞体脉冲、胞体和输出端的钙信号以及纹状体中的浓度都表现为误差，而不是期望$V$。两种解释哪一种在所有输出上都正确，仍在讨论中。 & \cite{Kim2020} & 已测量 \\
11 & 向后看：多巴胺报告因果联系的度量~\eqref{eq:retro} & 小鼠，在该理论与误差~\eqref{eq:ctd}预测不同的实验中测量伏隔核的多巴胺释放：在若干实验中释放遵循前者。 & \cite{Jeong2022} & 假说 \\
12 & 学习中多巴胺的响应逐渐从奖励移到预示信号 & 小鼠，学习过程中腹侧纹状体内\nt{DOPA}神经元输出端的信号：爆发逐渐从奖励时刻移到预示信号时刻，正如按~\eqref{eq:ctd}学习所要求的。 & \cite{Amo2022} & 已测量 \\
13 & 在价值不变时，\nt{DOPA}神经元对奖励味道的替换有响应 & 大鼠，腹侧被盖区神经元记录：把果汁换成价值相等的另一种会引起响应，尽管误差~\eqref{eq:ctd}为零。 & \cite{Takahashi2017} & 已测量 \\
14 & 人为诱发的爆发能教会两个中性刺激之间的联系 & 大鼠，无奖励的两刺激预配对实验：在从第一个刺激过渡到第二个时光遗传激活\nt{DOPA}神经元会建立联系，抑制它们则妨碍学习。 & \cite{Sharpe2017} & 已测量 \\
15 & 按特征的误差向量~\eqref{eq:vectortd} & 特征预测误差理论；向量形式没有直接检验。 & \cite{Gardner2018} & 假说 \\
16 & 同一任务中不同\nt{DOPA}神经元携带不同的量 & 虚拟走廊中的小鼠，\nt{DOPA}神经元双光子成像：有的与奖励相关，有的与速度和转向相关，有的与预示信号和选择准确性相关。 & \cite{Engelhard2019} & 已测量 \\
17 & 线束代替导线（命题~\ref{prop:bundle}） & 第~2--16~行的汇总：每种分解都已分别测得；它们都是同一信号的组成部分，这一统一图景没有实验支持。 & --- & 假说 \\
\bottomrule
\end{longtable}}
